\chapter{Weighing, measuring, crossing, and information}

\begin{orient}
\textbf{What this chapter is for.} A large family of puzzles is decided before you design a strategy, by an information count. A balance with three outcomes, used three times, distinguishes at most $27$ situations, so a puzzle that asks you to identify one of $12$ coins \emph{and} its direction of error, which is $24$ situations, is telling you that the bound is nearly tight and that no weighing may be wasted. Reading that bound first turns an open-ended search into a design problem with a specification. The chapter covers the information bound and how to use it, pouring problems and their arithmetic, state-space search done honestly, comparison lower bounds, the distributed-strategy family in which a parity code carries information nobody individually has, and then the geometric side: area invariants, tilings, symmetry, and extremal configurations.
\end{orient}

\section{Counting the information first}

\tech{The information-theoretic bound}{medium}
\cue A puzzle that limits the number of tests, queries, weighings or questions, and asks whether a given number suffices.
\method Count the outcomes. If each test has $b$ possible results and you are allowed $t$ of them, the entire procedure can distinguish at most $b^{t}$ situations, so a problem with $N$ possible truths requires $t\geq\log_{b}N$. Use the bound in both directions: it proves impossibility when $N>b^{t}$, and when $N$ is just under $b^{t}$ it tells you the strategy must split the possibilities as evenly as the arithmetic allows, which is a strong design constraint.

\begin{worked}
Twelve coins, one counterfeit of unknown weight direction, a balance, three weighings. There are $24$ possible truths and $3^{3}=27$ outcomes, so three weighings might suffice and no weighing may be badly unbalanced. The first weighing must therefore split $24$ into three groups as near to eight each as possible, which forces four coins against four, leaving four aside: eight truths in each pan outcome and eight in the balanced outcome. Every subsequent weighing is chosen by the same rule, and the strategy essentially writes itself once the balance requirement is recognised.
\end{worked}

\begin{trap}
Counting the objects rather than the truths. Twelve coins is not twelve possibilities, because the direction of the error is part of what must be determined; getting this count wrong is what makes the puzzle look impossible or trivially easy.
\end{trap}

\begin{seealsob}Comparison and sorting lower bounds; base representation; the pigeonhole principle.\end{seealsob}

\tech{Comparison and sorting lower bounds}{medium}
\cue Any question of the form ``how many comparisons are needed to''.
\method A comparison has two outcomes, so $t$ comparisons distinguish at most $2^{t}$ orderings. Sorting $n$ items must distinguish $n!$ of them, giving $t\geq\log_{2}n!\approx n\log_{2}n-n\log_{2}\eu$, which is the classical bound. Finding just the maximum needs only $n-1$ comparisons, since every non-maximum must lose at least once and each comparison produces one loser.

\begin{worked}
Find both the largest and the second largest of $n$ items with fewest comparisons. A knockout tournament finds the largest in $n-1$ comparisons; the second largest must have lost directly to the largest, and only $\ceilb{\log_{2}n}$ items did, so a second tournament among those costs $\ceilb{\log_{2}n}-1$. Total $n+\ceilb{\log_{2}n}-2$, which matches the information-style lower bound for this problem.
\end{worked}

\begin{seealsob}The information-theoretic bound; double counting.\end{seealsob}

\section{Arithmetic and state spaces}

\tech{Pouring problems via greatest common divisor}{easy}
\cue Two or more unmarked containers of known capacities and a target volume.
\method Every reachable volume is an integer combination of the capacities, so with capacities $a$ and $b$ the reachable amounts are exactly the multiples of $\gcd(a,b)$ that do not exceed $\max(a,b)$. Bezout's identity gives the coefficients, and the sign of each coefficient tells you which jar is being filled and which emptied.

\begin{worked}
Jugs of $5$ and $3$ litres, target $4$. Since $\gcd(5,3)=1$ divides $4$, it is reachable, and $4=3\cdot3-5$ suggests filling the three-litre jug three times and emptying the five-litre jug once. Concretely: fill the three, pour into the five; fill the three again, top up the five (which takes two), leaving one in the three-litre jug; empty the five, pour the one across, fill the three and add it, giving four.
\end{worked}

\begin{trap}
Assuming a target below the smaller capacity is easier. Reachability depends only on divisibility by the greatest common divisor, not on size, and $4$ from $\{5,3\}$ is no harder than $2$.
\end{trap}

\begin{seealsob}Base representation; explicit state-space search.\end{seealsob}

\tech{Explicit state-space search}{easy}
\cue Crossings, jumps, sliding pieces, or any puzzle with a small set of configurations and a legality rule on transitions.
\method Define the state as the smallest description that determines the future, list the legal transitions, and search breadth-first, which finds a shortest solution and proves it is shortest. Prune with an invariant if one exists: an invariant that separates the start from the goal turns the search into a one-line impossibility proof.

\begin{worked}
Three couples must cross a river in a boat holding two, with the constraint that no person may be with someone else's partner unless their own is present. The state is the triple of counts on the near bank plus the boat's position. Breadth-first search over the legal states finds a shortest solution in eleven crossings, and the same search with four couples and a two-seat boat finds none, which is the honest way to establish that the four-couple version is unsolvable without a larger boat.
\end{worked}

\begin{seealsob}Pouring problems; parity invariant; branch and contradict.\end{seealsob}

\tech{Distributed strategies and correlated errors}{hard}
\cue Several participants who cannot communicate after the puzzle begins, each seeing part of the situation, who must guess simultaneously. The question asks for a strategy guaranteeing a good outcome, and the naive analysis says the guesses are independent so nothing better than chance is possible.
\method The naive analysis is wrong because a strategy may correlate the errors. Nobody's individual guess beats chance, and that is not the objective: the objective is to concentrate all the wrong guesses onto a few situations so that the remaining situations are entirely right. A parity code is the standard device, since it lets each participant compute what the others' parity implies about their own unseen value.

\begin{worked}
Each of $n$ people receives a hat, black or white, chosen by independent fair coins, and sees everyone else's but not their own. All guess simultaneously; the team wins if everyone is correct. Without coordination the chance is $2^{-n}$. Agree in advance that everybody assumes the total number of black hats is even; then each person deduces their own hat from the parity of what they see. This wins exactly when the true parity is even, that is with probability $1/2$, for every $n$. All the losses have been concentrated on the odd-parity configurations, and each individual is still right exactly half the time.
\end{worked}

\begin{trap}
Concluding from ``each guess is individually a coin flip'' that the team cannot do better. Individual marginals say nothing about the joint distribution, and the whole technique lives in that gap.
\end{trap}

\begin{seealsob}Linearity of expectation; parity invariant; the information bound.\end{seealsob}

\section{Geometry and dissection}

\tech{Area and lattice invariants}{medium}
\cue A dissection, a rearrangement, or a claim about a polygon with vertices at lattice points.
\method Area is invariant under cutting and reassembling, which settles most rearrangement puzzles immediately, including the ones designed to look as though area has changed. For lattice polygons, Pick's theorem
\[
A=I+\frac{B}{2}-1,
\]
with $I$ interior and $B$ boundary lattice points, converts an area statement into a counting statement and back. A standard consequence: no equilateral triangle has all three vertices on the integer lattice, since Pick forces a rational area while the equilateral formula forces an irrational multiple of a square length.

\begin{trap}
Reading a picture instead of computing. The classical missing-square dissections work entirely because two slopes that look equal are not, and only arithmetic detects it.
\end{trap}

\begin{seealsob}Tiling impossibility by colouring; symmetry reduction.\end{seealsob}

\tech{Tiling impossibility by colouring}{medium}
\cue A region and a tile shape, and the question whether the region can be covered exactly.
\method Colour so that each placement of the tile covers a fixed multiset of colours, then count. This is the same technique as the colouring invariants of the invariants chapter, and it is listed again here because in geometric dress the colouring is usually driven by the tile's dimensions: a $1\times k$ tile wants colouring by index modulo $k$, an L-tromino wants a colouring that breaks the $2\times2$ symmetry, and a T-tetromino wants the chessboard.

\begin{worked}
Can $25$ T-tetrominoes tile a $10\times10$ board? Chequer the board, which then has $50$ squares of each colour. Every T covers three squares of one colour and one of the other. Suppose $k$ of the tiles cover three black squares and the remaining $25-k$ cover three white. Counting black squares gives $3k+(25-k)=2k+25=50$, so $2k=25$, which has no integer solution. No tiling exists.
\end{worked}

\begin{seealsob}Colouring arguments; area and lattice invariants.\end{seealsob}

\tech{Symmetry reduction and extremal configurations}{hard}
\cue A geometric optimisation, a claim that some configuration is best possible, or a counting problem over configurations that are equivalent under a symmetry group.
\method For optimisation, argue that an optimal configuration may be assumed symmetric, usually by a smoothing argument: replace an asymmetric pair by its symmetrised version and show the objective does not worsen. For counting, quotient by the symmetry group using the averaging formula. For existence and impossibility, take the extremal object and derive a contradiction from a local improvement.

\begin{worked}
Among all triangles inscribed in a given circle, the equilateral has the largest area. Fix two vertices; the third should maximise its distance from that chord, so it sits at the midpoint of the arc, making the triangle isosceles on that base. Applying the same argument to each of the three sides in turn forces all three to be equal. The smoothing step never decreases the area, so the equilateral triangle is optimal.
\end{worked}

\begin{trap}
Assuming the optimum is symmetric without a smoothing argument. It is true remarkably often and false often enough to matter, and the argument that establishes it is usually two lines.
\end{trap}

\begin{seealsob}The extremal principle; Burnside and symmetry counting.\end{seealsob}

\section{Recognition matrix}

\begin{recog}{@{}p{0.31\textwidth}p{0.35\textwidth}p{0.28\textwidth}@{}}
\rowcolor{mist}\mh{If you see} & \mh{Reach for} & \mh{If that fails} \\
\midrule
\endhead
A limited number of tests & Count truths against $b^{t}$ & Design for an even split \\
A balance, three outcomes & Powers of three; balanced ternary & Count directions of error too \\
``How many comparisons'' & $\log_{2}$ of the number of orderings & Tournament construction \\
Unmarked jugs, a target volume & Divisibility by the gcd; Bezout & Explicit state search \\
Crossings, jumps, sliding pieces & Breadth-first over states & An invariant to prove impossibility \\
Simultaneous guesses, no talking & A parity code; correlate the errors & Concentrate losses on few cases \\
``Each guess is a coin flip'' & Marginals do not bound the joint & Design the correlation \\
A dissection that gains area & Compute; the slopes differ & Pick's theorem \\
Lattice polygon, area asked & Pick: $A=I+B/2-1$ & Shoelace formula \\
A region and one tile shape & Colour to match the tile & Weighted colouring \\
Optimal configuration sought & Smoothing toward symmetry & Extremal principle \\
Configurations up to symmetry & Burnside averaging & Canonical form \\
\end{recog}
